\documentclass[12pt, letterpaper]{article}
\usepackage{parskip}
\setcounter{secnumdepth}{0}
\title{Statement Elaborato SWE}
\author{Francesco Giuntini}
\date{Gennaio 2026}

\begin{document}
\maketitle
\textbf{Tipologia elaborato:} Modo 2.\newline\newline
Si vuole realizzare un applicativo software per la \textbf{gestione
  dei programmi di allenamento} degli atleti di un personal trainer,
fornendo al PT la possibilità di creare ed assegnare schede
di allenamento.
\section{Descrizione}
Ci sono 2 tipi di utenti, ciascuno con diversi permessi:
\begin{itemize}
  \item \textbf{PT:} può visualizzare i dettagli dei suoi
    atleti, creare, assegnare e modificare le schede di tali
    atleti, accettare richieste di sessioni di allenamento da parte
    degli atleti e visualizzare tutte le sessioni di allenamento già prenotate.
  \item \textbf{Atleta:} può inserire e modificare i propri dati
    anagrafici, visualizzare le schede di allenamento a lui
    assegnate, richiedere sessioni di allenamento con il PT e
    visualizzare le proprie sessioni già prenotate.
\end{itemize}
Ciascun PT deve registrarsi per poter usare l'applicativo,
sarà il PT a registrare i propri atleti.\newline\newline
Una \textbf{scheda }raggruppa un numero di mesi di allenamento,
ciascun mese contiene i varia allenamenti composti a loro volta da
esercizi.\newline Una scheda può essere assegnata ad uno o più
atleti, e non si può assegnare ad un atleta due schede i cui
intervalli di tempo si sovrappongono (per esempio una scheda da
gennaio a marzo e una da febbraio ad aprile).\newline\newline
Un \textbf{esercizio} è definito dal suo nome, il gruppo muscolare
che allena, il peso da sollevare e il numero di serie e ripetizioni
da effettuare.\newline La tipologia di esercizi in una scheda varia
solo da mese a mese, ma il peso o le serie/ripetizioni possono
variare da allenamento a allenamento.\newline\newline
Una \textbf{richiesta di sessione di allenamento }con il PT
deve contenere un riferimento all'atleta che la richiede, data e ora
della sessione e la motivazione della sessione (ad esempio dubbi
sull'esecuzione degli esercizi).\newline\newline
Si prevede l'utilizzo di un \textbf{database} per la gestione dei
dati dell'applicativo.
\end{document}
